\section{Mask-sensitive bounds and limits of the certificate methods}
\label{sec:mask-methods}

Throughout this section $C$ is a linear source of dimension $D$ over an
arbitrary field, entries have dimension $b$, and the fixed shared decoders
may be nonlinear. Write $N=bd^2$. Probabilities sample receiver labels,
never inputs to a machine or native algorithms. A \emph{certificate} below
is the numerical right side of a specified upper-bound inequality. Lower
bounds on these numbers are not lower bounds on achievable code dimension.

\subsection{Finite and structured masks}

For a binary mask $p$ let $R_r(s)$ count starts of monochromatic runs of
length at least $r$ along addition by $s$: a start $c$ satisfies
$p(c-s)\ne p(c)=p(c+s)=\cdots=p(c+(r-1)s)$. Constant translation orbits
have no such starts. Put $Q_s=R_1(s)$, the number of directed color changes.

\begin{proposition}[A finite alternative]\label{prop:finite-two-thirds}
\formal{CapacityMaskFinite.lean}{finite_weighted_run_budget,IsSharedCode.four_mul_finrank_le,IsSharedCode.six_mul_finrank_le}
For a binary partition and $d>8$,
\[
 6D\le4bd^2+5bd.
\]
\end{proposition}

\begin{outline}
The two-stage nested inequality with shifts $(1,3)$ gives
$4D\le2N+bd(R_1+R_4+3)$. With shifts $(2,1)$ it gives
$4D\le2N+bd(2R_2+R_3+4)$: use $Q_2=2R_2$, and retain the possible
two translation orbits of the step $2$. The run counts here have step $1$.
Two copies of the first inequality plus one of the second use the budget
\[
 2R_1+2R_2+R_3+2R_4\le2d.
\]
For completeness, charge the two length-one tickets of a run to offset
$0$ in lanes $0,1$, its two length-two tickets to offset $1$ in those
lanes, its length-three ticket to offset $2$ in lane $0$, and its two
length-four tickets to offsets $2,3$ in lanes $1,0$. Tickets occur only
when the charged offset is within that run. Distinct runs cannot collide,
and tickets within one run have distinct position--lane pairs. Thus they
inject into two copies of the cycle. Substitution gives
$12D\le8N+10bd$.
\end{outline}

\begin{theorem}[An explicit family of masks]\label{thm:structured-mask}
\formal{CapacityMaskStructured.lean}{CapacityMaskStructured.first_component_count,CapacityMaskStructured.second_component_count,CapacityMaskStructured.structured_finrank_le}
Let $d=8\ell$, $\ell\ge1$, write $c=a+\ell r$ with $0\le a<\ell$ and
$r\in\mathbb Z/8$, and let
\[
 p(c)=v(a)\mathbin{\mathrm{xor}}u(r),\qquad
 u=(0,1,1,0,1,0,0,1),
\]
where $v$ is arbitrary. Then
\[
 D\le\frac9{16}bd^2+\frac12bd.
\]
\end{theorem}

\begin{outline}
Use nested shifts $3\ell$ and $1$. The first graph has exactly $2\ell$
components. For each fixed $a$, its cyclic order of high indices is
$0,3,6,1,4,7,2,5$, whose labels are $0,0,0,1,1,1,1,0$ up to a common
complement. Each color therefore induces one four-vertex path.

The second shift set is $\{0,3\ell,1,3\ell+1\}$. Its differences retain
the first graph and include $3\ell+1$. For $a<\ell-1$, that step sends
$(a,r)$ to $(a+1,r+3)$. For every $z,\delta\in\{0,1\}$ there is an $r$
with $u(r)=z$ and $u(r+3)=z\mathbin{\mathrm{xor}}\delta$: use
$r=0,6,1,2$ for $(z,\delta)=(0,0),(0,1),(1,0),(1,1)$, respectively.
Taking $z$ to be the desired global color xor $v(a)$ and
$\delta=v(a)\mathbin{\mathrm{xor}}v(a+1)$ connects that color's two
neighboring low-index components. All low indices are thereby connected
within each color. No edge joins colors, so the final graph has exactly
two components. When $\ell=1$ this was already true of the first graph.

The nested bound is $D\le N/2+n_1/4+n_2/4$, with
$n_1=bd^2/4$ and $n_2=2bd$. Substitution proves the claim. In particular,
binary digit parity at $d=2^h$, $h\ge3$, belongs to this family: $u$ is
the parity of the three high bits and $v$ the parity of the remaining bits.
The formal proof constructs classifiers of the actual graph quotients,
using the displayed paths and bridges to prove their fibers are exactly
the connected components.
\end{outline}

\begin{proposition}[Exact direct envelope for digit parity]
\label{prop:digit-parity-envelope}
\formal{CapacityDigitParityEnvelope.lean}{CapacityDigitParity.digit_parity_envelope,CapacityDigitParity.digit_parity_floor_envelope}
For cyclic binary digit parity at $d=2^h$, $h\ge3$, the minimum over all
nonempty shift sets $\Delta$ of the original direct translate certificate
is exactly $5bd^2/8$. This is the numerical envelope of that formula, not
the capacity of the mask.
\end{proposition}

\begin{outline}
Put $s_h(c)=(-1)^{p(c)}$ and
$\eta_h(k)=d^{-1}\sum_c s_h(c)s_h(c+k)$. The formal mask is the literal
parity of the lowest $h$ binary digits, defined by repeated division by
two. Its periodicity identifies the finite correlation with the cyclic
mask statistic. Splitting the average into even and odd positions gives
\[
 \eta_h(2k)=\eta_{h-1}(k),\qquad
 \eta_h(2k+1)=-\tfrac12\bigl(\eta_{h-1}(k)+\eta_{h-1}(k+1)\bigr).
\]
Every correlation lies in $[-1,1]$, since its summands are signs. The
recurrences show that each adjacent sum at level $h\ge1$ is either
$(\eta_{h-1}(k)-\eta_{h-1}(k+1))/2$ or its negative. It therefore lies in
$[-1,1]$. Every odd-index correlation at level $h\ge2$ lies in
$[-1/2,1/2]$. Even indices reduce to the preceding level, and the
nonzero correlations at $h=2$ are $0,-1,0$. Induction gives
$\eta_h(t)\le1/2$ for every nonzero cyclic shift. Consequently
\[
 Q_t=\tfrac d2(1-\eta_h(t))\ge d/4.
\]
The actual first-graph inequality $\chi(\{0,t\})\ge Q_t$ suffices for
the lower bound. Translating a shift set leaves its link relation
unchanged, so every two-element set has certificate at least $5N/8$.
For $q=|\Delta|\ge3$, the ambient term alone gives $2N/3$; a singleton
set gives $N$.

For attainment, set $\ell=d/8$ and split the three high bits from the
remaining bits. The literal parity mask is exactly the preceding
structured family, with $u$ the three-bit parity word. Its actual first
graph for $t=3\ell$ has $2\ell=d/4$ components. This two-element shift
set therefore attains $5N/8$. The value is integral, so flooring any
other certificate cannot improve it. The structured nested bound improves
this entire envelope by $bd(d-8)/16$ when $d>8$.
\end{outline}

\subsection{Run spectra and random masks}

We use the bounded-differences inequality in its following precise form.
For independent coordinates, changing coordinate $i$ changes a real function
$f$ by at most $c_i$; then either tail at distance $t\ge0$ from its mean
has probability at most $\exp(-2t^2/\sum_i c_i^2)$ when $\sum_i c_i^2>0$.
This is McDiarmid's inequality\cite{McDiarmid1989}; the displayed form and
the lower-tail reduction by $-f$ are stated in
\cite[Theorem~1, equation~(2)]{Warnke2012}. No independence between the
individual summands of $f$ is required. For the finite uniform product laws
used here, the shared \texttt{lean-machine} modules
\texttt{FiniteUniformProbability}, \texttt{FiniteBoundedDifferences}, and
\texttt{FiniteUniformWindows} prove the needed inequality and restriction
identities. The exponential-moment proof applies the bounded-interval
estimate successively to actual finite conditional averages; it retains
the full interval length $c_i$ at each coordinate.

\begin{proposition}[The three-stage spectrum]\label{prop:run-spectrum-three}
\formal{CapacityMaskRunSpectrum.lean}{IsSharedCode.run_spectrum_three}
For $d>1$, putting $R_r=R_r(1)$ gives
\[
 D\le \frac N2+bd\left(\frac{R_1}{4}+\frac{R_4}{8}
                                  +\frac{R_{13}}8\right)+\frac{3bd}4.
\]
\end{proposition}

\begin{outline}
Use the actual shifts $1,3,9$ in the three-stage inequality. Their
successive difference sets cover radii $1,4,13$. The corresponding
component counts are at most $R_1+1,R_4+2,R_{13}+2$, respectively;
the first step has a single translation orbit. Their weights are
$1/4,1/8,1/8$, so the additive error is $3bd/4$.
\end{outline}

\begin{theorem}[Almost every binary mask]\label{thm:random-mask-spectrum}
\formal{CapacityMaskRunProbability.lean}{CapacityMaskProbability.spectrum_bad_probability,CapacityMaskProbability.SpectrumGood.source_bound}
Let $d>13$, $\varepsilon\ge0$, and sample the $d$ mask bits independently
and fairly. Except on probability at most
$\exp(-128\varepsilon^2d/529)$, simultaneously every valid source and
decoder family for the sampled mask satisfy
\[
 D\le\left(\frac{41473}{65536}+\varepsilon\right)bd^2+\frac{3bd}4.
\]
\end{theorem}

\begin{outline}
A run-start indicator of length $r<d$ examines $r+1$ distinct bits, with
exactly two favorable assignments. Thus $\mathbb E R_r=d/2^r$.
Changing one bit affects at most $r+1$ of these indicators. For
$F=R_1/4+R_4/8+R_{13}/8$, the coordinate-change constant is therefore
$2/4+5/8+14/8=23/8$, while
$1/2+\mathbb EF/d=41473/65536$. Apply bounded differences with
$t=\varepsilon d$ and then the preceding deterministic inequality.
The event depends only on the mask, so there is no union over fields,
widths, codes or decoders. Conditioning on both colors occurring divides
the failure bound by $1-2^{1-d}$; identifying complementary labelings
does not change any statistic.
\end{outline}

The same reasoning gives the full fixed-depth spectrum. For $k\ge2$,
$d>(3^k-1)/2$, set $r_j=(3^j-1)/2$ and use shifts $3^{j-1}$. Balanced-ternary
digits show by induction that their difference set covers every integer
from $-r_j$ to $r_j$: the preceding interval and its translates by
$\pm3^{j-1}$ are consecutive and exhaustive. Each difference digit is
the difference of two digits in $\{0,1\}$. For a color in cyclic order,
a gap greater than $r$ corresponds to an intervening opposite-color run
of length at least $r$; its components number the larger of one and its
number of long gaps. Summing both colors gives $\chi_r\le R_r+2$;
at radius one the sharper $R_1+1$ also covers constant masks. Consequently
\[
 D\le\frac N2+bd\left(\sum_{j=1}^{k-1}\frac{R_{r_j}}{2^{j+1}}
                  +\frac{R_{r_k}}{2^k}\right)+\frac{3bd}4.
\]
Replace each run count by its mean to obtain the constant
$c_k=1/2+\sum_{j=1}^{k-1}2^{-r_j-j-1}+2^{-r_k-k}$. The coordinate-change
constant is $L_k=\sum_{j=1}^{k-1}(r_j+1)/2^{j+1}+(r_k+1)/2^k$, so for
$\varepsilon\ge0$ the upper tail at $\varepsilon d$ is at most
$\exp(-2\varepsilon^2d/L_k^2)$. The formal endpoints are
\lean{CapacityTernarySpectrum.Good.source_bound} for the source inequality
and \lean{CapacityTernarySpectrum.bad_probability} for its probability bound.
Their proofs compute graph counts, means and sensitivities on the actual
recursively generated translate sets.

\subsection{What the numerical methods cannot certify}

For nonnegative ambient sizes $n_j$ define the exact nested certificate
\[
 B_k(n)=\sum_{j=0}^{k-1}\frac{n_j}{2^{j+1}}+\frac{n_k}{2^k},
 \qquad B_0(n)=n_0.
\]
If $n_1\le n_0$, every $k\ge0$ satisfies
$B_k(n)\ge n_0/2+n_1/4$. For $k\ge2$ these are the first two terms;
for $k=1$ the first ambient has weight $1/2$; for $k=0$ use $n_1\le n_0$.

\begin{theorem}[The nested-component floor]\label{thm:component-method-floor}
\formal{CapacityMaskComparisonProbability.lean}{CapacityMaskProbability.component_floor_probability,CapacityMaskProbability.ChangeGood.certificate_floor}
For $\eta>0$, outside probability
$(d-1)\exp(-\eta^2d/2)$ under independent fair binary labels, every
nested component certificate starting from $n_0=N$ is at least
$(5/8-\eta/4)N$. Shifts and chain lengths may depend on the whole mask.
Taking $\eta=\sqrt{6\log d/d}$ gives failure probability less than $d^{-2}$
for $d>1$.
\end{theorem}

\begin{outline}
For every nonzero $s$, $\mathbb E Q_s=d/2$. A bit change affects at most
the two directed comparisons at that bit and its predecessor, so the
coordinate-change constant is $2$, even for a step of order two.
Bounded differences and a union over the $d-1$ nonzero steps give one
event on which all $Q_s\ge(1/2-\eta)d$. Every first graph has
$\chi(\{0,s\})\ge Q_s$: on a nonconstant translation cycle its
monochromatic components equal its directed changes; a constant cycle
adds a component and no change. The formal proof obtains the same
inequality by counting the actual retained step edges. Thus
$n_1\ge(1/2-\eta)N$; for a zero shift $n_1=N$ instead. Apply the first-term
bound for every chain. Finally
$(d-1)e^{-\eta^2d/2}=(d-1)d^{-3}<d^{-2}$ for the stated $\eta$.
\end{outline}

\begin{theorem}[The full-ambient isolation floor]\label{thm:isolation-method-floor}
\formal{CapacityMaskMethodEndpoints.lean}{CapacityMaskProbability.isolation_rate_probability,CapacityMaskProbability.IsolationGood.rounded_floor}
For $\eta\ge0$, start the isolation recurrence from \emph{all} $d-2$ nonrepresentative
rows of a nonempty binary partition. Write $n_j$ for the number of these
rows isolated in the graph of the $j$th nested shift set, in row units,
and use precisely $B_k(n)$ with this terminal ambient. Outside probability
$(d-1)\exp(-2\eta^2d/9)$ under independent fair labels, every resulting
completed numerical bound $D_{\rm cert}=bd(2+B_k(n))$ satisfies
\[
 \frac{D_{\rm cert}}{N}\ge\frac9{16}-\frac\eta4+\frac1{2d}.
\]
Final integer flooring preserves the safe bound
$9/16-\eta/4-1/(2d)$ for every $b\ge1$. With
$\eta=\sqrt{(27/2)\log d/d}$ the unconditional failure probability is
less than $d^{-2}$. Conditioning on both classes being nonempty divides
it by $1-2^{1-d}$.
\end{theorem}

\begin{outline}
For $s\ne0$, a vertex is isolated in $G_{\{0,s\}}$ precisely when its
label differs from both neighbors at $\pm s$. If $s$ has order two the
neighbors coincide and its probability is $1/2$; otherwise three distinct
fair bits give probability $1/4$. Hence the total $I_s$ has mean at least
$d/4$. A bit change affects at most the indicators centered at that bit
and its two neighbors, so its coordinate-change constant is $3$.
Bounded differences and a union over shifts give, simultaneously,
$I_s\ge d/4-\eta d$. Deleting any two representatives loses at most two
isolated vertices. Thus $n_1\ge d/4-\eta d-2$, independently of how the
representatives were selected. The zero shift has $n_1=d-2$ and also
satisfies this inequality. Since $n_0=d-2$ and all later sizes are
nonnegative, the first-term bound gives
$B_k\ge9d/16-\eta d/4-3/2$. Add the baseline $2$, multiply by $bd$, and
divide by $N$. Final integer flooring of $D_{\rm cert}$ reduces its normalized value
by less than $1/N\le1/d$. The stated logarithmic choice gives
$(d-1)e^{-2\eta^2d/9}=(d-1)d^{-3}<d^{-2}$.
\end{outline}

Both floors are preserved under minima, infima over nonempty families,
convex averages, and pointwise limits of these same certificates. Replacing
counts by upper estimates only raises their right sides. Neither theorem
covers a smaller actual starting support, a sharper terminal estimate,
additional compatibility equations, or signed combinations of inequalities.

\begin{theorem}[The direct method with many classes]\label{thm:direct-method-floor}
\formal{CapacityMaskComparisonProbability.lean}{CapacityMaskProbability.direct_floor_probability,CapacityMaskProbability.EqualityGood.direct_gain_inverse_sqrt}
Let $d\ge2$, $1\le M\le d$, and give the $d$ positions independent uniform
labels in $\{1,\ldots,M\}$. For a nonempty shift set $\Delta$, put
\[
 g_\Delta=\frac{d-\chi(\Delta)}{|\Delta|d},\qquad
 g_{\rm direct}=\max_{\varnothing\ne\Delta\subseteq\mathbb Z/d}g_\Delta.
\]
For $\eta>0$, except on probability at most
$(d-1)e^{-\eta^2d/2}$,
\[
 g_{\rm direct}\le\sqrt{\frac{1/M+\eta}{2}},\qquad
 \min_\Delta N(1-g_\Delta)
       \ge N\left(1-\sqrt{\frac{1/M+\eta}{2}}\right).
\]
In particular, for $\eta_d=\sqrt{6\log d/d}$, the failure probability is
less than $d^{-2}$. If $M\eta_d\le1$, then $g_{\rm direct}\le1/\sqrt M$.
\end{theorem}

\begin{outline}
For each nonzero $s$, let $A_s=|\{c:p(c+s)=p(c)\}|$. The endpoints
are distinct, so $\mathbb EA_s=d/M$. A coordinate change affects at most
two comparisons, even for an order-two step. Bounded differences and a
union over nonzero steps give one event on which
$A_s\le d(1/M+\eta)$ for all of them.

Fix any mask on this event, then any shift set of size $q$. Choose one
orientation for each of its $\binom q2$ unordered distinct pairs. All
graph edges generated by that pair number at most $A_s$ for the chosen
difference $s$. An order-two step double-counts an edge and repeated
differences may count it again; both only increase this upper bound.
Every graph component admits a spanning tree, so
\[
 d-\chi(\Delta)\le |E(G_\Delta)|
       \le\binom q2 d(1/M+\eta).
\]
The formal proof gives the same deficit inequality by the rank of the
edge-incidence map: its kernel is exactly the functions constant on the
actual translate-link quotient. Put $a=1/M+\eta$. Dividing and also using
$d-\chi\le d$ gives
$g_\Delta\le\min((q-1)a/2,1/q)\le\sqrt{a/2}$. The last step follows
from $g_\Delta^2\le(q a/2)g_\Delta\le a/2$; the gain is nonnegative.
This controls all sets from one event, even when chosen after inspecting
the labels. Substitute $\eta_d$, and use $a\le2/M$ when $M\eta_d\le1$.
\end{outline}

An absent label must not be counted as a nonempty class. If $B$ is the event
that every label occurs, the union bound gives
$\Pr(B^c)\le M(1-1/M)^d\le M e^{-d/M}=:q_{\rm empty}$ (at $M=1$ the
first value is zero). Thus the last theorem and $B$ hold simultaneously
with probability at least $1-d^{-2}-q_{\rm empty}$. When
$q_{\rm empty}<1$, its conditional failure probability is at most
$d^{-2}/(1-q_{\rm empty})$; otherwise use the exact positive $\Pr(B)$
as denominator. In the growing range $M\eta_d\le1$, the empty-label bound
tends to zero, since $M\le\sqrt{d/(6\log d)}$ and hence
$M e^{-d/M}\le\sqrt d\,e^{-\sqrt{6d\log d}}$ eventually. Every unlabeled
nonempty partition has exactly $M!$ labelings, so the same conditional
fraction holds for uniform unlabeled partitions. The formal quotient uses
equality of label fibers; it constructs the unique relabeling permutation
and proves the uniform-fiber identity. The conditioning statements and
the growing-range limit are formalized in
\texttt{CapacityMaskUnlabeled.lean} and \texttt{CapacityMaskGrowth.lean}. Together with the
universal difference-cover guarantee, this makes the order $M^{-1/2}$
sharp for that direct certificate method in the stated growing range.
It is not an optimal constant, an attained code rate, or a restriction on
the stronger nested methods. Final flooring changes a normalized direct
certificate by less than $1/N$.

As another comparison, putting $\eta=1/6$ in the binary change-count event
makes every direct two-shift certificate at least $2N/3$, outside
probability $(d-1)e^{-d/72}$. Every direct set of at least three shifts
already has that lower bound from its ambient term. For any fixed
$0<\varepsilon<2/3-41473/65536$ and sufficiently large $d$, the run-spectrum
bound is strictly below this entire direct envelope, except on probability
at most $e^{-128\varepsilon^2d/529}+(d-1)e^{-d/72}$. For rounded bounds
choose $d$ large enough that the difference exceeds one dimension.
No comparison here extracts a shared code from a native computation.
